\section{Quá trình Rà soát, Khắc phục và Kiểm chứng Hệ thống}

\subsection{Hành trình Khắc phục và Hoàn thiện (Remediation Journey)}
Một trong những giá trị cốt lõi của công trình này là tính minh bạch và sự nghiêm cẩn trong phương pháp luận khoa học. Trong quá trình phát triển hệ thống, nhóm nghiên cứu đã chủ động thực hiện ba đợt rà soát và khắc phục kỹ thuật toàn diện (được đánh mã hiệu P0, P1 và P2) nhằm loại bỏ các sai số đo lường và bảo vệ tính toàn vẹn của bằng chứng thực nghiệm \src{7}\src{8}.

\subsubsection{Giai đoạn P0: Thiết lập Quy chuẩn Báo cáo Bất biến}
Ở các phiên bản đầu, quy trình sinh báo cáo còn phụ thuộc vào các tệp kịch bản (\textit{scripts}) ad-hoc có khả năng ghi đè trực tiếp lên thư mục kết quả đánh giá cũ. Điều này tiềm ẩn nguy cơ làm biến đổi dữ liệu thực nghiệm đã chạy trước đó.

Nhóm đã tái cấu trúc toàn bộ quy trình báo cáo theo nguyên tắc **Dữ chứng Đóng băng (Frozen Evidence Pipeline)**:
\begin{itemize}
    \item Mọi kịch bản báo cáo chỉ được phép đọc dữ liệu từ các gói kết quả đã lưu và đã kiểm tra mã băm, tuyệt đối không được phép chỉnh sửa hay sinh lại số liệu trong quá trình tạo bảng biểu.
    \item Nếu phát hiện bất kỳ tệp dữ liệu nào bị thiếu hoặc sai lệch mã băm kiểm tra (\textit{checksum mismatch}), quy trình sẽ lập tức dừng lại (\textit{fail-closed}) và báo lỗi, từ chối tạo báo cáo trên dữ liệu không hoàn chỉnh.
\end{itemize}

\subsubsection{Giai đoạn P1: Loại bỏ Số liệu Mô phỏng và Đối soát Toàn vẹn Bằng chứng}
Trong đợt rà soát P1, nhóm phát hiện hai vấn đề nghiêm trọng về dữ liệu lịch sử trong kho mã nguồn:
\begin{enumerate}
    \item **Số liệu phân tích đánh đổi cũ được dựng từ hằng số:** Một số tài liệu phân tích đánh đổi (\textit{trade-off report}) trước đây sử dụng các bảng số đo được viết sẵn trong mã chương trình mà không qua chạy thực nghiệm thật với mô hình ngôn ngữ.
    \item **Nhãn con người được mô phỏng tự động:** Một tập 40 bộ nhãn hiệu chuẩn con người (\textit{human calibration}) trong quá khứ thực chất được sinh tự động bởi script thay vì do người thật chấm điểm.
\end{enumerate}

Với tinh thần trung thực học thuật tuyệt đối, nhóm đã **loại bỏ hoàn toàn các số liệu mô phỏng này khỏi toàn bộ báo cáo và kết luận khoa học** \src{8}. Các tài liệu kỹ thuật cũ đều được gắn cảnh báo đính chính rõ ràng. Nhóm đã xây dựng lại bộ công cụ chấm mù chuẩn hóa dành cho người thật (với phiếu chấm để trống) và chỉ công nhận kết quả khi có nhãn độc lập.

Đồng thời, nhóm tiến hành kiểm toán toàn diện tính toàn vẹn của gói kết quả gốc P8:
\begin{itemize}
    \item Bản sao lưu trong thư mục \texttt{output/.../p8-publication/} có dung lượng đầy đủ $10{,}6$ MB, chứa trọn vẹn tệp \texttt{preparation.json} và khớp hoàn toàn mã băm SHA-256 công bố (\texttt{35996008...997f}).
    \item Toàn bộ 720 lượt chạy thực nghiệm P8 được chứng minh là dữ liệu thật $100\%$, phản ánh chính xác hành vi của mô hình mà không có sự can thiệp nhân tạo.
\end{itemize}

\subsubsection{Giai đoạn P2: Cải tiến Chất lượng Truy xuất và Khắc phục Rơi Dữ kiện}
Dựa trên phân tích các trường hợp thất bại của hệ thống trong các đợt chạy thử nghiệm, nhóm đã xác định và khắc phục 6 nguyên nhân chính gây thất thoát thông tin trong pipeline tư vấn \src{7}:
\begin{enumerate}
    \item \textbf{Bảo toàn truy vấn tự nhiên gốc:} Thay vì rút gọn câu hỏi của người dùng thành các từ khóa rời rạc (làm mất các sắc thái phủ định hoặc điều kiện ràng buộc phụ), hệ thống giữ nguyên văn bản gốc của câu hỏi khi chuyển tiếp cho chuyên gia lập kế hoạch.
    \item \textbf{Cấp quyền kiểm tra tiêu đề đầy đủ khi chấm Lexical:} Khi đối sánh tên sách, hệ thống sử dụng tiêu đề ấn bản đã được cấp quyền thay vì chỉ so khớp chuỗi ký tự thô của từ khóa tìm kiếm.
    \item \textbf{Bảo toàn thuộc tính xuất xứ (Cross-Source Attribution):} Đảm bảo mọi dữ kiện được trích xuất từ đánh giá người mua luôn gắn chặt với mã sản phẩm tương ứng, không bị trôi ngữ cảnh khi ghép với thông tin danh mục.
    \item \textbf{Tối ưu hóa dung lượng trích dẫn theo ngân sách (Budget-Aware Snippet Expansion):} Tự động mở rộng độ dài đoạn trích dẫn văn bản nguồn khi ngân sách token của lượt hội thoại còn dư dả, giúp câu trả lời có thêm ngữ cảnh chi tiết mà không vi phạm hạn mức chi phí.
    \item \textbf{Bắt buộc kiểm tra đầy đủ các trường lọc:} Rà soát logic truy vấn cơ sở dữ liệu để đảm bảo không bỏ sót bất kỳ điều kiện lọc nào mà người dùng đã nêu (thể loại, khoảng giá, số trang, đánh giá tối thiểu).
    \item \textbf{Sửa đổi quy trình chạy kiểm thử không ghi đè:} Đảm bảo các runner kiểm định hoạt động độc lập và không làm biến đổi trạng thái của kho dữ liệu tĩnh.
\end{enumerate}

\subsection{Hệ thống Kiểm thử Hồi quy Tự động (Automated Regression Testing)}
Bên cạnh các đợt đánh giá chất lượng câu trả lời bằng mô hình ngôn ngữ lớn, Evidence Atlas được bảo vệ bởi một hệ thống kiểm thử tự động toàn diện gồm **1.456 ca kiểm thử** (\textit{automated unit and integration tests}) không gọi API bên ngoài, đảm bảo tính ổn định tuyệt đối của phần mềm trước mọi thay đổi mã nguồn:

\begin{table}[htbp]\small
\caption{Cấu trúc hệ thống kiểm thử tự động của Evidence Atlas.}
\begin{tabularx}{\linewidth}{lrrY}\toprule
Nhóm kiểm thử & Số lượng test & Thời gian thực thi & Phạm vi kiểm tra trọng tâm\\\midrule
\texttt{Security \& Gateway} & 218 & $\approx 12$ giây & Xác thực khóa API, phân quyền vai trò, ngăn chặn leo thang đặc quyền\\
\texttt{Contracts \& Schema} & 342 & $\approx 8$ giây & Tính bất biến của Pydantic models, kiểm tra định dạng giá, ID hợp lệ\\
\texttt{Supervisor Logic} & 415 & $\approx 25$ giây & Quy trình lập kế hoạch có chặn trên, giới hạn số bước, xử lý timeout\\
\texttt{Grounding \& Verification} & 265 & $\approx 15$ giây & Khớp xâu ký tự chính xác (\textit{exact-span}), định dạng trích dẫn \texttt{[C1]}\\
\texttt{PostgreSQL Integration} & 216 & $\approx 45$ giây & Khóa lease SQL, giao dịch cơ sở dữ liệu, kho tri thức vector bền vững\\\midrule
\textbf{Tổng cộng} & \textbf{1.456} & \textbf{$\approx 105$ giây} & \textbf{Tỷ lệ vượt qua: 100\% (0 lỗi)}\\\bottomrule
\end{tabularx}
\end{table}

Toàn bộ 1.456 ca kiểm thử trên đều chạy thành công trên mã nguồn hiện tại của dự án. Điều này chứng minh rằng mặc dù chất lượng sinh văn bản của mô hình ngôn ngữ lớn có thể dao động tùy thuộc vào độ phức tạp của câu hỏi, nhưng nền tảng phần mềm bên dưới—bao gồm cổng an toàn, cơ sở dữ liệu bền vững, bộ kiểm tra dữ chứng và cơ chế quản lý ngân sách—đều vận hành ổn định và đạt chuẩn kỹ nghệ phần mềm chuyên nghiệp.
